% !TEX program = lualatex
\documentclass[a4paper,10pt,twocolumn]{article}
\usepackage{kaitoushu}
\renewcommand{\baselinestretch}{1.2}
\lhead{応用数学　練習問題（採点用）}
\problemrange{第2章 §1 練習問題1（p.\,61）}

\begin{document}
\thispagestyle{fancy}

\begin{center}
  {\Large \textbf{応用数学　第2章 §1　練習問題1（p.\,61）}}
\end{center}
\vspace{0.5em}

\mondai{1}{
次の関数のラプラス変換を求めよ．\\
(1) $(t+1)^2$ \quad
(2) $\sin t \cos 2t$ \quad
(3) $(e^t + 3e^{-2t})\sin t$ \quad
(4) $t^2 \cos t$
}
\kotae{
(1) $(t+1)^2 = t^2+2t+1$ より
$\mathcal{L}[(t+1)^2] = \dfrac{2}{s^3}+\dfrac{2}{s^2}+\dfrac{1}{s} \quad (s>0)$\\[2pt]
(2) 積和公式：$\sin t\cos 2t = \dfrac{1}{2}(\sin 3t - \sin t)$ より
$\mathcal{L}[\sin t\cos 2t] = \dfrac{1}{2}\!\left(\dfrac{3}{s^2+9}-\dfrac{1}{s^2+1}\right) = \dfrac{s^2-3}{(s^2+1)(s^2+9)} \quad (s>0)$\\[2pt]
(3) 像関数の移動法則より
$\mathcal{L}[(e^t+3e^{-2t})\sin t] = \dfrac{1}{(s-1)^2+1}+\dfrac{3}{(s+2)^2+1} \quad (s>1)$\\[2pt]
(4) $\mathcal{L}[t^2 f(t)] = F''(s)$（$F(s)=\mathcal{L}[\cos t]=\dfrac{s}{s^2+1}$）を用いる．

$F'(s) = \dfrac{1-s^2}{(s^2+1)^2}$，\quad
$F''(s) = \dfrac{2s(s^2-3)}{(s^2+1)^3}$

$\mathcal{L}[t^2\cos t] = \dfrac{2s(s^2-3)}{(s^2+1)^3} \quad (s>0)$
}

\mondai{2}{
正の定数 $a, b, k, l$ について，次の関数を単位ステップ関数で表し，
そのラプラス変換を求めよ．\\
(1) $f(t) = \begin{cases} k & (a \leq t < b) \\ 0 & (t < a,\ t \geq b) \end{cases}$
\quad
(2) $g(t) = \begin{cases} 0 & (t < 0) \\ k & (0 \leq t < a,\ t \geq b) \\ l & (a \leq t < b) \end{cases}$
}
\kotae{
(1) $f(t) = k\bigl[U(t-a) - U(t-b)\bigr]$ より
\[
  \mathcal{L}[f(t)] = \dfrac{k(e^{-as} - e^{-bs})}{s} \quad (s>0)
\]
(2) $g(t) = kU(t) + (l-k)U(t-a) + (k-l)U(t-b)$ より
\[
  \mathcal{L}[g(t)] = \dfrac{k + (l-k)e^{-as} + (k-l)e^{-bs}}{s} \quad (s>0)
\]
}

\mondai{3}{
次の関数の逆ラプラス変換を求めよ．\\
(1) $\dfrac{s}{s^2+4}$ \quad
(2) $\dfrac{s+1}{s^2+9}$ \quad
(3) $\dfrac{1}{4s^2-1}$ \quad
(4) $\dfrac{s}{(s+1)^2}$ \quad
(5) $\dfrac{s+3}{(s+1)(s-2)}$ \quad
(6) $\dfrac{s+1}{(2s-1)^2}$
}
\kotae{
(1) $\mathcal{L}^{-1}\!\left[\dfrac{s}{s^2+2^2}\right] = \cos 2t$\\[2pt]
(2) $\dfrac{s+1}{s^2+9} = \dfrac{s}{s^2+9}+\dfrac{1}{s^2+9}$ より
$\mathcal{L}^{-1} = \cos 3t + \dfrac{1}{3}\sin 3t$\\[2pt]
(3) $4s^2-1=(2s-1)(2s+1)$ を部分分数分解：
$\dfrac{1}{(2s-1)(2s+1)} = \dfrac{1/4}{s-1/2} - \dfrac{1/4}{s+1/2}$
\[
  \mathcal{L}^{-1} = \dfrac{1}{4}e^{t/2} - \dfrac{1}{4}e^{-t/2}
\]
(4) $s=(s+1)-1$ より $\dfrac{s}{(s+1)^2} = \dfrac{1}{s+1}-\dfrac{1}{(s+1)^2}$
\[
  \mathcal{L}^{-1} = e^{-t} - te^{-t} = (1-t)e^{-t}
\]
(5) $\dfrac{s+3}{(s+1)(s-2)} = \dfrac{A}{s+1}+\dfrac{B}{s-2}$ とおく．\\
通分：$s+3=A(s-2)+B(s+1)$\\
$s=-1$：$2=-3A \Rightarrow A=-\dfrac{2}{3}$，\quad $s=2$：$5=3B \Rightarrow B=\dfrac{5}{3}$
\[
  \mathcal{L}^{-1} = -\dfrac{2}{3}e^{-t}+\dfrac{5}{3}e^{2t}
\]
(6) $s+1=\dfrac{1}{2}(2s-1)+\dfrac{3}{2}$ より
$\dfrac{s+1}{(2s-1)^2} = \dfrac{1/2}{2s-1}+\dfrac{3/2}{(2s-1)^2} = \dfrac{1/4}{s-1/2}+\dfrac{3/8}{(s-1/2)^2}$
\[
  \mathcal{L}^{-1} = \dfrac{1}{4}e^{t/2}+\dfrac{3}{8}te^{t/2} = \dfrac{(2+3t)e^{t/2}}{8}
\]
}

\mondai{4}{
関数 $F(s) = \dfrac{s^2}{s^3+8}$ について，次の問いに答えよ．\\
(1) 次の形式が成り立つように $A, B, C$ の値を定めよ：
$F(s) = \dfrac{A}{s+2} + \dfrac{Bs+C}{s^2-2s+4}$\\
(2) 関数 $F(s)$ の逆ラプラス変換を求めよ．
}
\kotae{
(1) $s^3+8=(s+2)(s^2-2s+4)$ より，両辺に $(s+2)(s^2-2s+4)$ を掛けて
$s^2 = A(s^2-2s+4) + (Bs+C)(s+2)$．\\
$s=-2$ で $4 = 12A \Rightarrow A = \dfrac{1}{3}$．\\
$s^2$ の係数: $1 = A + B \Rightarrow B = \dfrac{2}{3}$．\\
定数項: $0 = 4A + 2C \Rightarrow C = -2A = -\dfrac{2}{3}$．
\[
  A = \dfrac{1}{3},\quad B = \dfrac{2}{3},\quad C = -\dfrac{2}{3}
\]
(2) $\dfrac{Bs+C}{s^2-2s+4} = \dfrac{2}{3}\cdot\dfrac{s-1}{(s-1)^2+3}$ より
\[
  \mathcal{L}^{-1}[F(s)] = \dfrac{1}{3}e^{-2t} + \dfrac{2}{3}e^{t}\cos\sqrt{3}\,t
\]
}

\mondai{5}{
原関数の移動法則を用いて，次の関数の逆ラプラス変換を求めよ．\\
(1) $\dfrac{e^{-s}}{s^3}$ \quad
(2) $\dfrac{e^{-2\pi s}}{s^2+1}$
}
\kotae{
(1) $\mathcal{L}^{-1}\!\left[\dfrac{1}{s^3}\right] = \dfrac{t^2}{2}$ と原関数の移動法則 $\mathcal{L}^{-1}[e^{-as}F(s)] = f(t-a)U(t-a)$ より
\[
  \mathcal{L}^{-1}\!\left[\dfrac{e^{-s}}{s^3}\right] = \dfrac{(t-1)^2}{2}U(t-1)
\]
(2) $\mathcal{L}^{-1}\!\left[\dfrac{1}{s^2+1}\right] = \sin t$ と原関数の移動法則より
\[
  \mathcal{L}^{-1}\!\left[\dfrac{e^{-2\pi s}}{s^2+1}\right] = \sin(t-2\pi)\,U(t-2\pi) = \sin t \cdot U(t-2\pi)
\]
}

\end{document}
